\documentclass[aspectratio=169]{beamer}
\usepackage[utf8]{inputenc}
\usepackage[spanish]{babel}
\usepackage{amsmath}
\usepackage{graphicx}
\usepackage{booktabs}
\usepackage{tikz}

% Tema minimalista y estético
\usetheme{Boadilla}
\usecolortheme{whale}

% Paleta de colores personalizados (Estilo Moderno/Tech)
\definecolor{DarkBlueTech}{RGB}{15, 32, 67}
\definecolor{AccentCyan}{RGB}{0, 173, 181}
\definecolor{LightGray}{RGB}{240, 243, 244}

\setbeamercolor{structure}{fg=DarkBlueTech}
\setbeamercolor{palette primary}{bg=DarkBlueTech, fg=white}
\setbeamercolor{palette secondary}{bg=AccentCyan, fg=white}
\setbeamercolor{palette tertiary}{bg=DarkBlueTech, fg=white}
\setbeamercolor{title}{bg=DarkBlueTech, fg=white}
\setbeamercolor{item}{fg=AccentCyan}

% Quitar los botones de navegación de beamer
\setbeamertemplate{navigation symbols}{}

% Formato de título
\title[Recomendador Híbrido]{\textbf{Sistema Híbrido de Recomendación de Laptops}}
\subtitle{Arquitectura en Cascada de 2 Niveles para E-Commerce de Hardware}
\author[Roger Vilcacari]{Roger Vilcacari}
\institute[UNAP]{Sistemas de Recomendación - Proyecto Final}
\date{\today}

\begin{document}

% -------------------------
% Diapositiva 1
% -------------------------
\begin{frame}
    \titlepage
\end{frame}

% -------------------------
% Diapositiva 2
% -------------------------
\begin{frame}{Introducción y Motivación}
    \begin{itemize}
        \item \textbf{Contexto:} La elección de una laptop es un proceso complejo. Involucra alta inversión económica y la necesidad de emparejar requisitos técnicos estrictos con perfiles de uso muy variados (estudiantes, programadores, arquitectos, gamers).
        \vspace{0.3cm}
        \item \textbf{Problema Central:} Los sistemas de recomendación actuales en plataformas de E-Commerce suelen basarse únicamente en popularidad o historial de compras, ignorando la incompatibilidad técnica (Ej. recomendar una laptop Celeron a un diseñador 3D solo porque es "muy vendida").
        \vspace{0.3cm}
        \item \textbf{Objetivo:} Desarrollar un sistema experto híbrido que garantice el 100\% de cumplimiento de restricciones técnicas y que además ofrezca una experiencia personalizada mitigando la falta de historial.
    \end{itemize}
\end{frame}

% -------------------------
% Diapositiva 3
% -------------------------
\begin{frame}{Definición de los Retos Técnicos}
    \begin{block}{Reto 1: El \textit{Cold Start} (Inicio en frío)}
        Las personas renuevan sus laptops cada 3-5 años. En la mayoría de sesiones, el usuario es "nuevo" para el sistema. El \textbf{Filtrado Colaborativo} puro colapsa sin datos de interacción previos.
    \end{block}
    
    \vspace{0.2cm}
    \begin{block}{Reto 2: Restricciones Duras (Constraints)}
        En dominios de entretenimiento (cine, música), un falso positivo genera disgusto temporal. En hardware, un falso positivo (ej. recomendar una laptop sin GPU para tareas pesadas) invalida toda la compra.
    \end{block}

    \vspace{0.2cm}
    \begin{block}{Reto 3: Espacio Multidimensional}
        La relevancia no se define por "marca", sino por un vector compuesto por RAM, Almacenamiento, Tasa de Refresco, GPU, CPU, etc.
    \end{block}
\end{frame}

% -------------------------
% Diapositiva 4
% -------------------------
\begin{frame}{El Dataset de Trabajo}
    Para el entorno de validación, se procesó información de escenarios reales:
    
    \vspace{0.3cm}
    \begin{itemize}
        \item \textbf{Fuente de Datos:} Dataset "Laptop Price Prediction" de Kaggle (2024).
        \item \textbf{Volumen:} 1,273 modelos únicos de laptops reales en el mercado.
        \item \textbf{Preprocesamiento y Extracción:}
        \begin{itemize}
            \item Conversión automatizada de monedas a USD.
            \item Extracción mediante \textit{Regex} de componentes vitales: Cantidad de RAM (GB), núcleos de CPU, existencia de GPU Dedicada y soporte CUDA.
            \item Creación de un "Screen Score" basado en resoluciones y tecnologías de panel.
        \end{itemize}
        \item \textbf{Simulación Colaborativa:} Generación de casi 3,000 interacciones (ratings) sintéticas divididas en arquetipos (Gamers, Estudiantes, Editores) para entrenar el modelo latente.
    \end{itemize}
\end{frame}

% -------------------------
% Diapositiva 5
% -------------------------
\begin{frame}{Solución: Arquitectura Híbrida en Cascada}
    \begin{center}
    \resizebox{!}{0.7\textheight}{
    \begin{tikzpicture}[
        box/.style={draw=DarkBlueTech, thick, rounded corners, align=center, fill=LightGray, text width=6cm, minimum height=1.2cm},
        arrow/.style={->, thick, DarkBlueTech, >=stealth}
    ]
        % Nodos
        \node[box] (input) at (0, 3.5) {\textbf{Input del Usuario}\\Presupuesto + Perfil Cualitativo};
        
        \node[box, fill=DarkBlueTech!10] (nivel1) at (0, 1.75) {\textbf{Nivel 1: Motor Constraint-Based}\\Aplica reglas de exclusión lógicas};
        
        \node[box, fill=AccentCyan!10, text width=9cm] (nivel2) at (0, 0) {\textbf{Nivel 2: Ensamble Dinámico Ponderado}\\$U_{maut}$ (Knowledge) + $S_{content}$ (Content) + $\hat{r}_{svd}$ (Collaborative)};
        
        \node[box] (output) at (0, -1.75) {\textbf{Ranking Final (Top-K Recomendaciones)}};

        % Flechas
        \draw[arrow] (input) -- (nivel1);
        \draw[arrow] (nivel1) -- node[right, font=\footnotesize] {Solo Pasan ítems válidos} (nivel2);
        \draw[arrow] (nivel2) -- (output);
    \end{tikzpicture}
    }
    \end{center}
\end{frame}

% -------------------------
% Diapositiva 6
% -------------------------
\begin{frame}{Nivel 1: Inferencia Lógica (Filtro Duro)}
    Actúa como la barrera principal del sistema. Transforma las especificaciones del usuario en \textbf{reglas binarias (Constraints)}.
    
    \vspace{0.3cm}
    \begin{block}{Construcción Matemática de la Máscara}
        Para cada laptop $i$, se genera un escalar $M_{rules}$:
        \begin{equation}
            M_{rules}(i) = \begin{cases} 
            1 & \text{si } P_i \le B \land (R_i \ge R_{min}) \land (V_i \ge V_{min}) \\
            0 & \text{en cualquier otro caso}
            \end{cases}
        \end{equation}
    \end{block}

    \vspace{0.3cm}
    \begin{itemize}
        \item \textbf{Ejemplo:} Si el perfil es \textit{Data Science}, se fuerza $RAM \ge 16GB$ y GPU con CUDA.
        \item Cualquier equipo con $M_{rules}(i) = 0$ queda descartado permanentemente, reduciendo el espacio de búsqueda y eliminando la fricción de falsos positivos.
    \end{itemize}
\end{frame}

% -------------------------
% Diapositiva 7
% -------------------------
\begin{frame}{Nivel 2 (Sub-Modelo A): Teoría de Utilidad (MAUT)}
    El \textit{Multi-Attribute Utility Theory} (Knowledge-Based) permite puntuar laptops incluso si \textbf{ningún usuario las ha comprado antes}.
    
    \vspace{0.2cm}
    El usuario define la importancia (pesos $w$) que le da a distintas dimensiones. El modelo normaliza (Min-Max) los valores técnicos de la laptop ($S$).
    
    \begin{block}{Cálculo de Utilidad $U_{maut}$}
        \begin{equation}
            U_{maut}(i) = w_{precio} \cdot S_{p}(i) + w_{rendimiento} \cdot S_{r}(i) + w_{pantalla} \cdot S_{s}(i) + w_{portab.} \cdot S_{m}(i)
        \end{equation}
    \end{block}
    
    \begin{itemize}
        \item Las métricas se normalizan al rango $[0, 1]$.
        \item Ej: Para $S_{precio}$, los equipos más económicos obtienen mayor puntaje cercano a $1$.
    \end{itemize}
\end{frame}

% -------------------------
% Diapositiva 8
% -------------------------
\begin{frame}{Nivel 2 (Sub-Modelo B): Factorización SVD}
    Modelo \textbf{Collaborative Filtering} diseñado para descubrir patrones de consumo de la comunidad que no son obvios matemáticamente (ej. lealtad de marca).
    
    \begin{block}{Predicción de Calificación $\hat{r}$}
        \begin{equation}
            \hat{r}_{ui} = \mu + b_u + b_i + q_i^T p_u
        \end{equation}
    \end{block}
    
    \begin{itemize}
        \item $\mu$: Promedio global.
        \item $b_u, b_i$: Sesgos del usuario (criticismo) y del ítem (popularidad global).
        \item $q_i^T p_u$: Producto punto de los vectores latentes (Factores Ocultos).
        \item Se entrena minimizando el error con \textbf{Stochastic Gradient Descent (SGD)}.
        \item \textit{Normalización:} El resultado (de 1 a 5 estrellas) se escala a $[0,1]$ para sumarse al híbrido.
    \end{itemize}
\end{frame}

% -------------------------
% Diapositiva 9
% -------------------------
\begin{frame}{Nivel 2 (Sub-Modelo C): Similitud Coseno}
    Modelo \textbf{Content-Based} que mide la proximidad geométrica entre lo que el usuario necesita y lo que la laptop ofrece.
    
    \begin{block}{Métrica de Distancia Vectorial}
        Se genera un vector ideal del usuario $v_u$ y un vector de la laptop $v_i$.
        \begin{equation}
            S_{content}(i) = \cos(v_u, v_i) = \frac{v_u \cdot v_i}{\|v_u\| \|v_i\|}
        \end{equation}
    \end{block}
    
    \begin{itemize}
        \item El vector incluye: \texttt{[RAM\_Score, CPU\_Score, GPU\_Score, Calidad\_Pantalla]}.
        \item Ignora el precio, centrándose exclusivamente en el paralelismo de las capacidades técnicas.
    \end{itemize}
\end{frame}

% -------------------------
% Diapositiva 10
% -------------------------
\begin{frame}{Ecuación Híbrida y Estrategia de Pesos Dinámicos}
    Los tres modelos se integran en una combinación lineal ponderada, condicionada por el Nivel 1.
    
    \begin{block}{Función de Puntaje Global}
    \begin{equation}
        Score_i = M_{rules}(i) \times \Big[ \alpha \cdot U_{maut}(i) + \beta \cdot \hat{r}_{svd\_norm}(i) + \gamma \cdot S_{content}(i) \Big]
    \end{equation}
    \end{block}
    
    \vspace{0.2cm}
    \textbf{Mitigación Inteligente del Cold Start:}
    Los coeficientes cambian según la madurez del usuario.
    \begin{itemize}
        \item \textbf{Usuario Nuevo (Cold Start):} $\alpha = 0.60$ (Fuerte peso a reglas y preferencias), $\gamma = 0.30$, $\beta = 0.10$ (Se apaga casi por completo la predicción comunitaria).
        \item \textbf{Usuario Frecuente:} $\alpha = 0.30, \beta = 0.50$ (Se le da poder a la sabiduría colaborativa).
    \end{itemize}
\end{frame}

% -------------------------
% Diapositiva 11
% -------------------------
\begin{frame}{Implementación del Sistema Web (Dashboard)}
    Para validar su funcionamiento, se construyó una interfaz E-Commerce funcional:
    
    \begin{itemize}
        \item \textbf{Backend:} Servidor en Python usando \texttt{FastAPI} de alto rendimiento.
        \item \textbf{Almacenamiento:} Modelos pre-entrenados cargados en memoria y base de datos vectorizada.
        \item \textbf{Frontend:} SPA interactiva (HTML5, JS, CSS Glassmorphism) que permite ajustar los factores de peso y probar distintos perfiles en tiempo real.
        \item \textbf{Endpoints:} \texttt{/api/recommend} recibe el JSON de preferencias y ejecuta la cascada en fracciones de segundo.
    \end{itemize}
\end{frame}

% -------------------------
% Diapositiva 12
% -------------------------
\begin{frame}{Evaluación Científica: Constraint Satisfaction Rate (CSR)}
    Se evaluó qué porcentaje de los Top-10 recomendados eran realmente \textbf{útiles} para el usuario (cumplían sus mínimos técnicos):
    
    \vspace{0.4cm}
    \begin{columns}
        \column{0.5\textwidth}
        \begin{table}
            \centering
            \begin{tabular}{lc}
                \toprule
                \textbf{Algoritmo} & \textbf{CSR@10} \\
                \midrule
                SVD Puro & 31.7\% \\
                Content-Based Puro & 48.3\% \\
                \textbf{Arquitectura Híbrida} & \textbf{93.3\%} \\
                \bottomrule
            \end{tabular}
        \end{table}
        
        \column{0.5\textwidth}
        \begin{block}{Análisis}
        El SVD tradicional recomienda ítems populares que a menudo violan el presupuesto o el uso específico. El Sistema Híbrido, gracias a $M_{rules}$, es virtualmente impecable.
        \end{block}
    \end{columns}
\end{frame}

% -------------------------
% Diapositiva 13
% -------------------------
\begin{frame}{Evaluación Científica: Precisión y RMSE}
    \begin{columns}
        \column{0.5\textwidth}
        \textbf{Predictividad Colaborativa (RMSE):}
        \begin{itemize}
            \item RMSE del modelo SVD en dataset de test: \textbf{1.033}
            \item Demuestra que el algoritmo SVD logra predecir el interés a 1 estrella de margen.
        \end{itemize}
        
        \column{0.5\textwidth}
        \textbf{Precisión@10 y Recall@10:}
        \begin{itemize}
            \item En escenarios de \textit{Cold Start}, las métricas clásicas son engañosamente bajas, pero la integración del vector de conocimiento MAUT asegura que las recomendaciones sean objetivamente relevantes para el perfil consultado.
        \end{itemize}
    \end{columns}
\end{frame}

% -------------------------
% Diapositiva 14
% -------------------------
\begin{frame}{Conclusiones y Trabajo Futuro}
    \begin{itemize}
        \item \textbf{Hibridación en Cascada:} Un modelo puramente basado en IA generativa o colaborativa no es seguro para comercio de hardware. Los sistemas basados en restricciones (\textit{Constraint-based}) son imperativos.
        \item \textbf{Solución MAUT:} Transformar specs técnicas crudas a una "función de utilidad" normalizada resolvió la asimetría de los datos y permitió recomendaciones perfectas sin interacciones pasadas.
        \item \textbf{Trabajo Futuro:} 
        \begin{itemize}
            \item Integrar procesamiento de lenguaje natural (NLP) para que el usuario escriba sus necesidades en texto libre.
            \item Migración a Bases de Datos Vectoriales (\textit{Milvus} o \textit{Pinecone}) para escalar a catálogos de millones de productos.
        \end{itemize}
    \end{itemize}
\end{frame}

% -------------------------
% Diapositiva 15
% -------------------------
\begin{frame}[plain]
    \begin{center}
        \vspace{1cm}
        \Huge \textbf{¡Muchas Gracias!} \\
        \vspace{0.5cm}
        \Large Arquitectura Híbrida de 2 Niveles para E-Commerce \\
        \vspace{1cm}
        \normalsize \textit{Turno de Preguntas y Respuestas}
    \end{center}
\end{frame}

\end{document}
